%!TeX root=../gardentop.tex
\chapter{The Robin who Showed the Way} 
	
\begin{figure}[t!]
\centering
\includegraphics[width=\linewidth]{july}
\end{figure}

 \lettrine[lines=6]{S}{he} looked at the key quite a long time. She turned it over and over, and thought about it. As I have said before, she was not a child who had been trained to ask permission or consult her elders about things. All she thought about the key was that if it was the key to the closed garden, and she could find out where the door was, she could perhaps open it and see what was inside the walls, and what had happened to the old rose-trees. It was because it had been shut up so long that she wanted to see it. It seemed as if it must be different from other places and that something strange must have happened to it during ten years. Besides that, if she liked it she could go into it every day and shut the door behind her, and she could make up some play of her own and play it quite alone, because nobody would ever know where she was, but would think the door was still locked and the key buried in the earth. The thought of that pleased her very much.

Living as it were, all by herself in a house with a hundred mysteriously closed rooms and having nothing whatever to do to amuse herself, had set her inactive brain to working and was actually awakening her imagination. There is no doubt that the fresh, strong, pure air from the moor had a great deal to do with it. Just as it had given her an appetite, and fighting with the wind had stirred her blood, so the same things had stirred her mind. In India she had always been too hot and languid and weak to care much about anything, but in this place she was beginning to care and to want to do new things. Already she felt less »contrary,« though she did not know why.

She put the key in her pocket and walked up and down her walk. No one but herself ever seemed to come there, so she could walk slowly and look at the wall, or, rather, at the ivy growing on it. The ivy was the baffling thing. Howsoever carefully she looked she could see nothing but thickly-growing, glossy, dark green leaves. She was very much disappointed. Something of her contrariness came back to her as she paced the walk and looked over it at the tree-tops inside. It seemed so silly, she said to herself, to be near it and not be able to get in. She took the key in her pocket when she went back to the house, and she made up her mind that she would always carry it with her when she went out, so that if she ever should find the hidden door she would be ready.

Mrs~Medlock had allowed Martha to sleep all night at the cottage, but she was back at her work in the morning with cheeks redder than ever and in the best of spirits.

»I got up at four o'clock,« she said. »Eh! it was pretty on th' moor with th' birds gettin' up an' th' rabbits scamperin' about an' th' sun risin'. I didn't walk all th' way. A man gave me a ride in his cart an' I can tell you I did enjoy myself.«

She was full of stories of the delights of her day out. Her mother had been glad to see her and they had got the baking and washing all out of the way. She had even made each of the children a dough-cake with a bit of brown sugar in it.

»I had 'em all pipin' hot when they came in from playin' on th' moor. An' th' cottage all smelt o' nice, clean hot bakin' an' there was a good fire, an' they just shouted for joy. Our Dickon he said our cottage was good enough for a king to live in.«

In the evening they had all sat round the fire, and Martha and her mother had sewed patches on torn clothes and mended stockings and Martha had told them about the little girl who had come from India and who had been waited on all her life by what Martha called »blacks« until she didn't know how to put on her own stockings.

»Eh! they did like to hear about you,« said Martha. »They wanted to know all about th' blacks an' about th' ship you came in. I couldn't tell 'em enough.«

Mary reflected a little.

»I'll tell you a great deal more before your next day out,« she said, »so that you will have more to talk about. I dare say they would like to hear about riding on elephants and camels, and about the officers going to hunt tigers.«

»My word!« cried delighted Martha. »It would set 'em clean off their heads. Would tha' really do that, Miss? It would be same as a wild beast show like we heard they had in York once.«

»India is quite different from Yorkshire,« Mary said slowly, as she thought the matter over. »I never thought of that. Did Dickon and your mother like to hear you talk about me?«

»Why, our Dickon's eyes nearly started out o' his head, they got that round,« answered Martha. »But mother, she was put out about your seemin' to be all by yourself like. She said, 'Hasn't Mr~Craven got no governess for her, nor no nurse?' and I said, 'No, he hasn't, though Mrs~Medlock says he will when he thinks of it, but she says he mayn't think of it for two or three years.'«

»I don't want a governess,« said Mary sharply.

»But mother says you ought to be learnin' your book by this time an' you ought to have a woman to look after you, an' she says: 'Now, Martha, you just think how you'd feel yourself, in a big place like that, wanderin' about all alone, an' no mother. You do your best to cheer her up,' she says, an' I said I would.«

Mary gave her a long, steady look.

»You do cheer me up,« she said. »I like to hear you talk.«

Presently Martha went out of the room and came back with something held in her hands under her apron.

»What does tha' think,« she said, with a cheerful grin. »I've brought thee a present.«

»A present!« exclaimed Mistress Mary. How could a cottage full of fourteen hungry people give any one a present!

»A man was drivin' across the moor peddlin',« Martha explained. »An' he stopped his cart at our door. He had pots an' pans an' odds an' ends, but mother had no money to buy anythin'. Just as he was goin' away our 'Lizabeth Ellen called out, 'Mother, he's got skippin'-ropes with red an' blue handles.' An' mother she calls out quite sudden, 'Here, stop, mister! How much are they?' An' he says 'Tuppence,' an' mother she began fumblin' in her pocket an' she says to me, 'Martha, tha's brought me thy wages like a good lass, an' I've got four places to put every penny, but I'm just goin' to take tuppence out of it to buy that child a skippin'-rope,' an' she bought one an' here it is.«

She brought it out from under her apron and exhibited it quite proudly. It was a strong, slender rope with a striped red and blue handle at each end, but Mary Lennox had never seen a skipping-rope before. She gazed at it with a mystified expression.

»What is it for?« she asked curiously.

»For!« cried out Martha. »Does tha' mean that they've not got skippin'-ropes in India, for all they've got elephants and tigers and camels! No wonder most of 'em's black. This is what it's for; just watch me.«

And she ran into the middle of the room and, taking a handle in each hand, began to skip, and skip, and skip, while Mary turned in her chair to stare at her, and the queer faces in the old portraits seemed to stare at her, too, and wonder what on earth this common little cottager had the impudence to be doing under their very noses. But Martha did not even see them. The interest and curiosity in Mistress Mary's face delighted her, and she went on skipping and counted as she skipped until she had reached a hundred.

»I could skip longer than that,« she said when she stopped. »I've skipped as much as five hundred when I was twelve, but I wasn't as fat then as I am now, an' I was in practice.«

Mary got up from her chair beginning to feel excited herself.

»It looks nice,« she said. »Your mother is a kind woman. Do you think I could ever skip like that?«

»You just try it,« urged Martha, handing her the skipping-rope. »You can't skip a hundred at first, but if you practise you'll mount up. That's what mother said. She says, 'Nothin' will do her more good than skippin' rope. It's th' sensiblest toy a child can have. Let her play out in th' fresh air skippin' an' it'll stretch her legs an' arms an' give her some strength in 'em.'«

It was plain that there was not a great deal of strength in Mistress Mary's arms and legs when she first began to skip. She was not very clever at it, but she liked it so much that she did not want to stop.

»Put on tha' things and run an' skip out o' doors,« said Martha. »Mother said I must tell you to keep out o' doors as much as you could, even when it rains a bit, so as tha' wrap up warm.«

Mary put on her coat and hat and took her skipping-rope over her arm. She opened the door to go out, and then suddenly thought of something and turned back rather slowly.

»Martha,« she said, »they were your wages. It was your twopence really. Thank you.« She said it stiffly because she was not used to thanking people or noticing that they did things for her. »Thank you,« she said, and held out her hand because she did not know what else to do.

Martha gave her hand a clumsy little shake, as if she was not accustomed to this sort of thing either. Then she laughed.

»Eh! tha' art a queer, old-womanish thing,« she said. »If tha'd been our 'Lizabeth Ellen tha'd have give me a kiss.«

Mary looked stiffer than ever.

»Do you want me to kiss you?«

Martha laughed again.

»Nay, not me,« she answered. »If tha' was different, p'raps tha'd want to thysel'. But tha' isn't. Run off outside an' play with thy rope.«

Mistress Mary felt a little awkward as she went out of the room. Yorkshire people seemed strange, and Martha was always rather a puzzle to her. At first she had disliked her very much, but now she did not.

The skipping-rope was a wonderful thing. She counted and skipped, and skipped and counted, until her cheeks were quite red, and she was more interested than she had ever been since she was born. The sun was shining and a little wind was blowing—not a rough wind, but one which came in delightful little gusts and brought a fresh scent of newly turned earth with it. She skipped round the fountain garden, and up one walk and down another. She skipped at last into the kitchen-garden and saw Ben Weatherstaff digging and talking to his robin, which was hopping about him. She skipped down the walk toward him and he lifted his head and looked at her with a curious expression. She had wondered if he would notice her. She really wanted him to see her skip.

»Well!« he exclaimed. »Upon my word! P'raps tha' art a young 'un, after all, an' p'raps tha's got child's blood in thy veins instead of sour buttermilk. Tha's skipped red into thy cheeks as sure as my name's Ben Weatherstaff. I wouldn't have believed tha' could do it.«

»I never skipped before,« Mary said. »I'm just beginning. I can only go up to twenty.«

»Tha' keep on,« said Ben. »Tha' shapes well enough at it for a young 'un that's lived with heathen. Just see how he's watchin' thee,« jerking his head toward the robin. »He followed after thee yesterday. He'll be at it again to-day. He'll be bound to find out what th' skippin'-rope is. He's never seen one. Eh!« shaking his head at the bird, »tha' curosity will be th' death of thee sometime if tha' doesn't look sharp.«

Mary skipped round all the gardens and round the orchard, resting every few minutes. At length she went to her own special walk and made up her mind to try if she could skip the whole length of it. It was a good long skip and she began slowly, but before she had gone half-way down the path she was so hot and breathless that she was obliged to stop. She did not mind much, because she had already counted up to thirty. She stopped with a little laugh of pleasure, and there, lo and behold, was the robin swaying on a long branch of ivy. He had followed her and he greeted her with a chirp. As Mary had skipped toward him she felt something heavy in her pocket strike against her at each jump, and when she saw the robin she laughed again.

»You showed me where the key was yesterday,« she said. »You ought to show me the door to-day; but I don't believe you know!«

The robin flew from his swinging spray of ivy on to the top of the wall and he opened his beak and sang a loud, lovely trill, merely to show off. Nothing in the world is quite as adorably lovely as a robin when he shows off—and they are nearly always doing it.

Mary Lennox had heard a great deal about Magic in her Ayah's stories, and she always said that what happened almost at that moment was Magic.

One of the nice little gusts of wind rushed down the walk, and it was a stronger one than the rest. It was strong enough to wave the branches of the trees, and it was more than strong enough to sway the trailing sprays of untrimmed ivy hanging from the wall. Mary had stepped close to the robin, and suddenly the gust of wind swung aside some loose ivy trails, and more suddenly still she jumped toward it and caught it in her hand. This she did because she had seen something under it—a round knob which had been covered by the leaves hanging over it. It was the knob of a door.

She put her hands under the leaves and began to pull and push them aside. Thick as the ivy hung, it nearly all was a loose and swinging curtain, though some had crept over wood and iron. Mary's heart began to thump and her hands to shake a little in her delight and excitement. The robin kept singing and twittering away and tilting his head on one side, as if he were as excited as she was. What was this under her hands which was square and made of iron and which her fingers found a hole in?

It was the lock of the door which had been closed ten years and she put her hand in her pocket, drew out the key and found it fitted the keyhole. She put the key in and turned it. It took two hands to do it, but it did turn.

And then she took a long breath and looked behind her up the long walk to see if any one was coming. No one was coming. No one ever did come, it seemed, and she took another long breath, because she could not help it, and she held back the swinging curtain of ivy and pushed back the door which opened slowly—slowly.

Then she slipped through it, and shut it behind her, and stood with her back against it, looking about her and breathing quite fast with excitement, and wonder, and delight.

She was standing inside the secret garden.